% ==============================================================================
% Section: The Allocative Cost of Price Controls
% Status: COMPLETE (2026-03-28 Session R4)
% ==============================================================================

\section{The Allocative Cost of Price Controls}
\label{sec:allocative_cost}

The Shapiro settlement imposes a hard cap at $\bar{p}_S = \$325$/MW-day UCAP
on PJM's capacity demand curve for the 2026/27 and 2027/28 delivery years.
Section~\ref{sec:results} established that the SFE equilibrium $p^* = \$329$
exceeds $\bar{p}_S$ in 2026/27 --- the settlement is binding by \$4/MW-day at
the RTO level and by \$115--\$141/MW-day in 2025/26 constrained LDAs. Binding
a price cap below the equilibrium changes market outcomes through two channels:
a revenue transfer from capacity sellers to buyers, and a distortion of the
entry and exit incentives that the capacity price is designed to convey. This
section analyzes both channels. Section~\ref{sec:21billion} quantifies the
revenue transfer in relation to the \$21~billion savings claim.

% ------------------------------------------------------------------------------
\subsection{Revenue Transfer and the Net CONE Benchmark}
\label{subsec:alloc_transfer}
% ------------------------------------------------------------------------------

Table~\ref{tab:cap_comparison} summarizes the key price levels and the implied
revenue transfer per MW-year under the settlement. For the 2026/27 delivery
year, the SFE equilibrium falls at the VRR Point~(a) price $\bar{p} = \$329.17$/MW-day,
yielding a binding gap of \$4.17/MW-day relative to the settlement cap. At
365 days per year, this corresponds to a revenue transfer of approximately
\$1,522 per MW-year at the RTO level. For 2027/28, the VRR Point~(a) price
under the new design is \$333.44/MW-day, implying a binding gap of \$8.44/MW-day
and a revenue transfer of approximately \$3,081 per MW-year.

\begin{table}[ht]
  \centering
  \caption{Settlement Price Levels vs.\ SFE Equilibrium and Net CONE (2026/27 and 2027/28)}
  \label{tab:cap_comparison}
  \begin{tabular}{lcccccc}
    \hline\hline
    Year & $\bar{p}$ & $\bar{p}_S$ & $p^*$ & Net CONE & $\bar{p}_f$ &
      $p^* - \bar{p}_S$ \\
    \hline
    2026/27 & \$329.17 & \$325.00 & \$329.17 & \$212.14 & \$175.00 & \$4.17 \\
    2027/28 & \$333.44 & \$325.00 & \$333.44 & \$242.52 & \$179.55 & \$8.44 \\
    \hline\hline
  \end{tabular}
  \smallskip

  \noindent\footnotesize
  \textit{Notes:} All prices in \$/MW-day UCAP. $\bar{p}$ is the VRR Point~(a)
  price cap; $\bar{p}_S = \$325$/MW-day is the Shapiro settlement cap;
  $p^* = \bar{p}$ because the SFE clears at the cap for $K \leq 3$ (both years
  observed at-cap in BRA data); Net CONE is the published Net Cost of New Entry
  for the RTO zone \citep{PJM2025_settlement_slides}; $\bar{p}_f$ is the
  settlement floor (Point~(b)/(c) in the new VRR curve); $p^* - \bar{p}_S$ is
  the binding gap. Revenue transfer per MW-year $= (p^* - \bar{p}_S) \times 365$.
\end{table}

The revenue transfer from sellers to buyers per MW-year is modest relative to
the settlement's stated savings --- \$1,522 and \$3,081 per MW-year in 2026/27
and 2027/28 respectively, compared to the \$21~billion aggregate claim, which
uses a counterfactual price of over \$500/MW-day and a much larger binding gap.
Section~\ref{sec:21billion} reconciles these estimates. The relevant comparison
here is not the size of the transfer per MW-year, but its proximity to the Net
CONE investment threshold. In 2026/27, Net CONE is \$212.14/MW-day. The
settlement cap at \$325/MW-day is well above Net CONE (ratio 1.53), so in
isolation the cap does not prevent new entry from being financially viable:
a generator earning \$325/MW-day exceeds its Net CONE by 53\%. The investment
distortion from the cap alone is therefore limited in 2026/27.

% ------------------------------------------------------------------------------
\subsection{Entry Distortion and the Informational Role of Prices}
\label{subsec:alloc_entry}
% ------------------------------------------------------------------------------

The more substantive allocative concern is not whether prices remain above
Net CONE today, but whether suppressing the price signal degrades the
forward-looking information that investors use to evaluate entry decisions
for the 2028/29 delivery year and beyond. \citet{Hayek1945_knowledge}
argues that prices perform an indispensable coordinative function: they
aggregate and transmit dispersed private knowledge about scarcity conditions
that no central authority can replicate. In a capacity market, the scarcity
signal is the gap between the clearing price and Net CONE. A price of
\$329/MW-day signals acute scarcity (\$117/MW-day above Net CONE in 2026/27);
capping it at \$325/MW-day preserves the directional signal but compresses
the magnitude.

The informational loss is larger when the cap is binding in consecutive
years. If the 2027/28 auction also clears at the settlement cap (\$325 with
$\bar{p} = \$333.44$ and Net CONE $= \$242.52$), potential entrants observe
only the cap rather than the unconstrained equilibrium price. They cannot
distinguish between ``the market would have cleared at \$333, signaling
\$91/MW-day above Net CONE'' and ``the market is capped and might have
cleared anywhere above \$325.'' This uncertainty reduces the informativeness
of the capacity price for investment decisions, a concern also raised in the
price-caps-in-SFE literature \citep{Anderson2005_sfe_pricecaps}.

The settlement's applicability to constrained LDAs amplifies this effect.
In 2026/27, all LDA-level SFE equilibria fall at or near the respective
VRR Point~(a) caps (Section~\ref{subsec:res_lda}), but the binding gap
between $p^*$ and $\bar{p}_S$ is structurally larger in tightly constrained
LDAs because those LDAs face steeper demand curves and higher Net CONE values
(e.g., PS and PS NORTH Net CONE $= \$369$/MW-day in 2026/27, far above both
$\bar{p}_S$ and $\bar{p}$). In such LDAs, the settlement cap truncates a
price signal that was already struggling to attract investment.

% ------------------------------------------------------------------------------
\subsection{The Floor's Distinct Allocative Role}
\label{subsec:alloc_floor}
% ------------------------------------------------------------------------------

The \$175/MW-day floor is not the mirror image of the cap. It operates as
an investment adequacy mechanism: all capacity offered at or below \$175/MW-day
UCAP is procured at the floor price, regardless of whether the market would
have cleared there \citep{PJM2025_settlement_slides}. PJM characterizes the
floor explicitly as ``supporting near-term investment'' --- it prevents the
price from collapsing in over-supplied market conditions, protecting the
revenue stream for marginal capacity.

A natural benchmark for evaluating the floor is Net CONE. If the floor exceeds
Net CONE, it provides a revenue guarantee above the investment threshold,
functioning as an implicit subsidy to existing capacity that may not be needed
for reliability. At the RTO level, the 2026/27 floor ($\$177.24$/MW-day in
VRR parameters) is \emph{below} Net CONE ($\$212.14$/MW-day): the floor does
not guarantee full cost recovery and thus does not over-compensate existing
generators on net. In the SWMAAC zone, however, the relationship inverts:
SWMAAC Net CONE is \$171.02/MW-day while the floor is \$177.24/MW-day, so the
floor exceeds the local investment threshold by \$6.22/MW-day and provides
above-cost revenue to SWMAAC capacity.

The floor's interaction with the cap creates a price band: outcomes between
\$175 and \$325 are interior equilibria; outcomes outside this band are
truncated. For $K = 3$, the SFE equilibrium in 2026/27 falls at \$329 ---
above the band. The settlement cap therefore does not allow the capacity
price to reach the strategic equilibrium; it instead pins the outcome to
the top of the band. From the perspective of the standard missing-money
literature \citep{Cramton2005_capacity, Joskow2008_capacity_payments}, a
price floor is a design-level intervention to correct the missing money
problem in energy-only markets. Applying a floor simultaneously with a
cap creates an internally inconsistent incentive: the cap signals that
scarcity is not severe enough to warrant prices above \$325, while the
floor signals that prices below \$175 are inadequate to support investment.
The two signals are not contradictions only if the competitive equilibrium
falls in the interior of the band --- which it does not under the $K = 3$
structural regime the SFE model identifies.

% ------------------------------------------------------------------------------
\subsection{The Price Band as a Non-Competitive Equilibrium}
\label{subsec:alloc_band}
% ------------------------------------------------------------------------------

\citet{GlaeserLuttmer2003_rent} demonstrate in the context of residential
rent control that binding price controls generate misallocation beyond the
standard deadweight triangle: when quantity is rationed below the market-
clearing level, allocation is distorted away from highest-willingness-to-pay
consumers. The direct mechanism --- relationship-based allocation in rental
housing --- does not translate to capacity markets, where all capacity
offering at or below the cap clears. The relevant analogue for capacity
markets is inter-temporal: a binding cap suppresses the price signal that
should direct investment toward future delivery years with the highest
expected scarcity, creating an efficiency loss that accumulates dynamically
rather than appearing as a static misallocation of current capacity.

The settlement band (\$175--\$325) defines a price interval that is
neither the competitive equilibrium ($c \approx \$150$, far below the
floor) nor the model-predicted benchmark price ($p^* \approx \$329$--\$333,
above the cap). It is an administratively imposed outcome. Whether this
intervention is welfare-improving depends on whether the gap between $p^*$
and $\bar{p}_S$ --- of order \$4--\$8/MW-day --- justifies the associated
information loss and investment distortion from cap-binding. The SFE model
suggests this gap is modest (1--3\% above the cap) at the RTO level;
the cap's revenue transfer benefit is small relative to its informational
cost under the model's assumptions. Section~\ref{sec:policy_alternatives}
considers whether structural remedies can reduce $p^*$ itself, rather than
truncating the equilibrium price signal.
